% <<<PDFPAGE 319>>>
\chapter{概率论与统计学常用结论}
\label{app:probability}
\bookimprint{https://doi.org/10.1007/978-3-031-64087-2}

\section{常用概率分布}
\subsection{伯努利分布}
\[
\begin{aligned}
\mathrm{B}(x\mid\theta)&=\theta^x(1-\theta)^{1-x},
  \qquad \text{其中 }x\in\{0,1\}\text{，}\theta\in[0,1]\\
\mathbb{E}[x]&=\theta\\
\operatorname{Var}[x]&=\theta(1-\theta)
\end{aligned}
\]

\subsection{类别分布（Multinoulli 分布）}
\[
\mathrm{M}(\mathbf{x}\mid\boldsymbol{\theta})=\prod_{d=1}^D\theta_d^{x_d},
\]
其中对所有 $d=1,2,\ldots,D$，有 $x_d\in\{0,1\}$、$\theta_d\in[0,1]$，且 $\sum_{d=1}^D\theta_d=1$。
\[
\begin{aligned}
\mathbb{E}[x_d]&=\theta_d\\
\operatorname{Var}[x_d]&=\theta_d(1-\theta_d)
\end{aligned}
\]
\emph{译注。} 这里的 $\mathbf{x}$ 是独热向量，恰有一个分量等于 1。

\subsection{正态分布}
\[
\begin{aligned}
\mathcal{N}(x\mid\mu,\sigma^2)
 &=\frac{1}{\sqrt{2\pi}\,\sigma}
   \exp\left\{-\frac{(x-\mu)^2}{2\sigma^2}\right\}\\
\mathbb{E}[x]&=\mu\\
\operatorname{Var}[x]&=\sigma^2
\end{aligned}
\]

% <<<PDFPAGE 320>>>
\subsection{多元正态分布}
\[
\mathcal{N}(\mathbf{x}\mid\boldsymbol{\mu},\boldsymbol{\Sigma})
=\frac{1}{(2\pi)^{D/2}}\frac{1}{|\boldsymbol{\Sigma}|^{1/2}}
 \exp\left\{-\tfrac12(\mathbf{x}-\boldsymbol{\mu})^{\mathsf T}
 \boldsymbol{\Sigma}^{-1}(\mathbf{x}-\boldsymbol{\mu})\right\},
\]
其中 $\mathbf{x}$ 为 $D$ 维向量，$\boldsymbol{\mu}$ 为 $D$ 维均值向量，$\boldsymbol{\Sigma}$ 为 $D\times D$ 的协方差矩阵。
\[
\begin{aligned}
\mathbb{E}[\mathbf{x}]&=\boldsymbol{\mu}\\
\operatorname{Cov}[\mathbf{x}]&=\boldsymbol{\Sigma}.
\end{aligned}
\]

\subsection{Beta 分布}
\[
\operatorname{Beta}(x\mid a,b)=\frac{\Gamma(a+b)}{\Gamma(a)\Gamma(b)}
                             x^{a-1}(1-x)^{b-1},
\]
其中 $x\in[0,1]$、$a>0$、$b>0$，且 $\Gamma(x)=\int_0^\infty t^{x-1}e^{-t}\,\mathrm{d}t$。
\[
\begin{aligned}
\mathbb{E}[x]&=\frac{a}{a+b}\\
\operatorname{Var}[x]&=\frac{ab}{(a+b)^2(a+b+1)}.
\end{aligned}
\]

\subsection{边缘分布}
连续情形：
\[
p(x)=\int p(x,y)\,\mathrm{d}y
\]
离散情形：
\[
p(x)=\sum_y p(x,y)
\]

\subsection{条件分布}
\[
p(y\mid x)=\frac{p(x,y)}{p(x)}
\]

\subsection{多元正态分布的边缘分布和条件分布}
假设 $\mathbf{x}\sim\mathcal{N}(\mathbf{x}\mid\boldsymbol{\mu},\boldsymbol{\Sigma})$，其中
% <<<PDFPAGE 321>>>
\[
\mathbf{x}=\begin{bmatrix}\mathbf{x}_a\\\mathbf{x}_b\end{bmatrix},
\qquad
\boldsymbol{\mu}=\begin{bmatrix}\boldsymbol{\mu}_a\\\boldsymbol{\mu}_b\end{bmatrix},
\qquad
\boldsymbol{\Sigma}
=\begin{bmatrix}\boldsymbol{\Sigma}_a&\boldsymbol{\Sigma}_c\\
                \boldsymbol{\Sigma}_c^{\mathsf T}&\boldsymbol{\Sigma}_b\end{bmatrix}.
\]
那么有以下关系：
\[
\begin{aligned}
p(\mathbf{x}_a)&=\mathcal{N}(\mathbf{x}_a\mid\boldsymbol{\mu}_a,\boldsymbol{\Sigma}_a),\\
p(\mathbf{x}_a\mid\mathbf{x}_b)
 &=\mathcal{N}(\mathbf{x}_a\mid\hat{\boldsymbol{\mu}}_a,\hat{\boldsymbol{\Sigma}}_a),
\end{aligned}
\]
其中
\[
\begin{aligned}
\hat{\boldsymbol{\mu}}_a
 &=\boldsymbol{\mu}_a+\boldsymbol{\Sigma}_c\boldsymbol{\Sigma}_b^{-1}
   (\mathbf{x}_b-\boldsymbol{\mu}_b),\\
\hat{\boldsymbol{\Sigma}}_a
 &=\boldsymbol{\Sigma}_a-\boldsymbol{\Sigma}_c\boldsymbol{\Sigma}_b^{-1}
                                   \boldsymbol{\Sigma}_c^{\mathsf T}.
\end{aligned}
\]

\subsection{求和法则}
\[
p(x)=\sum_y p(x,y)
\]

\subsection{乘积法则}
\[
\begin{aligned}
p(x,y)&=p(x\mid y)p(y)\\
      &=p(y\mid x)p(x)
\end{aligned}
\]

\subsection{贝叶斯法则}
\[
p(y\mid x)=\frac{p(x\mid y)p(y)}{p(x)}
\]

\section{统计学}
\subsection{最大似然估计量}
给定来自同一分布 $p(\mathbf{x}\mid\theta)$ 的 $N$ 个独立样本 $\mathcal{D}=\{\mathbf{x}_1,\ldots,\mathbf{x}_N\}$。似然函数为
\[
p(\mathcal{D}\mid\theta)=\prod_{n=1}^{N}p(\mathbf{x}_n\mid\theta).
\]
似然函数 $p(\mathcal{D}\mid\theta)$ 的对数为
\[
\log p(\mathcal{D}\mid\theta)=\sum_{n=1}^{N}\log p(\mathbf{x}_n\mid\theta).
\]
原文称参数的最大似然估计量（maximum likelihood estimator）$\theta_{\mathrm{ML}}$ 使似然函数最小。给出的式子为
% <<<PDFPAGE 322>>>
\[
p(\mathcal{D}\mid\theta_{\mathrm{ML}})=\max_\theta p(\mathcal{D}\mid\theta).
\]
\emph{译注。} 按原式，应是使似然函数最大。

\subsection{最大后验估计量}
给定来自同一分布 $p(\mathbf{x}\mid\theta)$ 的 $N$ 个独立样本 $\mathcal{D}=\{\mathbf{x}_1,\ldots,\mathbf{x}_N\}$。原文称参数的最大后验估计量（maximum a posteriori estimator，MAP）$\theta_{\mathrm{MAP}}$ 使后验分布最小。给出的式子为
\[
p(\theta_{\mathrm{MAP}}\mid\mathcal{D})=\max_\theta p(\theta\mid\mathcal{D}).
\]
\emph{译注。} 按原式，应是使后验分布最大。

\subsection{决策中的风险}
风险（risk）即期望损失，定义为
\[
\mathcal{R}[\bar{y}]=\iint L(y,\bar{y}(\mathbf{x}))\,p(\mathbf{x},y)\,
                         \mathrm{d}\mathbf{x}\,\mathrm{d}y,
\]
其中 $L(\cdot,\cdot)$ 是损失函数。
